\documentclass[12pt,a4paper]{article}
\usepackage[UTF8]{ctex}
\usepackage{amsmath,amssymb,amsthm}
\usepackage{geometry}
\usepackage{graphicx}
\usepackage{hyperref}
\usepackage{bm}

\geometry{margin=2.5cm}

\title{$[\text{ad}_{V_i}] A_i \dot{\theta}_i$与$[\text{ad}_{V_i}](A_i) \dot{\theta}_i$的区别与展开表示}
\author{Modern Robotics}
\date{}

\begin{document}

\maketitle

\section{引言}

在牛顿-欧拉逆向动力学的推导中，我们经常遇到两种看似相似但含义略有不同的表示：
\begin{itemize}
\item $[\text{ad}_{V_i}] A_i \dot{\theta}_i$
\item $[\text{ad}_{V_i}](A_i) \dot{\theta}_i$
\end{itemize}

本文档将详细解释这两种表示的区别，并给出它们的完整展开形式。

\section{符号定义}

\subsection{基本符号}

\begin{itemize}
\item $V_i = (\omega_i, v_i) \in \mathbb{R}^6$：连杆$i$的twist，其中$\omega_i \in \mathbb{R}^3$是角速度，$v_i \in \mathbb{R}^3$是线速度
\item $A_i = (\omega_{A_i}, v_{A_i}) \in \mathbb{R}^6$：关节$i$的螺旋轴
\item $\dot{\theta}_i \in \mathbb{R}$：关节$i$的速度（标量）
\item $[\text{ad}_{V_i}] \in \mathbb{R}^{6 \times 6}$：$\text{ad}_{V_i}$算子的矩阵表示
\end{itemize}

\subsection{$\text{ad}$算子的定义}

对于twist $V = (\omega, v)$，$\text{ad}_V$算子的矩阵表示为：
\begin{equation}
[\text{ad}_V] = \begin{bmatrix}
[\omega] & 0 \\
[v] & [\omega]
\end{bmatrix} \in \mathbb{R}^{6 \times 6}
\label{eq:ad_matrix}
\end{equation}

其中$[\omega]$和$[v]$分别是$\omega$和$v$的反对称矩阵表示。

\section{两种表示的区别}

\subsection{表示1：$[\text{ad}_{V_i}] A_i \dot{\theta}_i$}

\textbf{含义}：
\begin{itemize}
\item 这是\textbf{矩阵乘法}的表示
\item 计算顺序：先计算$[\text{ad}_{V_i}] A_i$（矩阵乘以向量），再乘以标量$\dot{\theta}_i$
\item 等价于：$( [\text{ad}_{V_i}] A_i ) \dot{\theta}_i$
\end{itemize}

\textbf{展开形式}：

设$V_i = (\omega_i, v_i)$，$A_i = (\omega_{A_i}, v_{A_i})$，则：

\begin{align}
[\text{ad}_{V_i}] A_i \dot{\theta}_i &= \begin{bmatrix}
[\omega_i] & 0 \\
[v_i] & [\omega_i]
\end{bmatrix} \begin{bmatrix}
\omega_{A_i} \\
v_{A_i}
\end{bmatrix} \dot{\theta}_i \\
&= \begin{bmatrix}
[\omega_i] \omega_{A_i} \\
[v_i] \omega_{A_i} + [\omega_i] v_{A_i}
\end{bmatrix} \dot{\theta}_i \\
&= \begin{bmatrix}
\omega_i \times \omega_{A_i} \\
v_i \times \omega_{A_i} + \omega_i \times v_{A_i}
\end{bmatrix} \dot{\theta}_i \\
&= \begin{bmatrix}
(\omega_i \times \omega_{A_i}) \dot{\theta}_i \\
(v_i \times \omega_{A_i} + \omega_i \times v_{A_i}) \dot{\theta}_i
\end{bmatrix}
\label{eq:notation1_expanded}
\end{align}

\subsection{表示2：$[\text{ad}_{V_i}](A_i) \dot{\theta}_i$}

\textbf{含义}：
\begin{itemize}
\item 这是\textbf{算子作用}的表示
\item $[\text{ad}_{V_i}](A_i)$表示$\text{ad}_{V_i}$算子作用在$A_i$上的结果
\item 根据定义：$\text{ad}_{V_i}(A_i) = [\text{ad}_{V_i}] A_i$
\item 因此：$[\text{ad}_{V_i}](A_i) = [\text{ad}_{V_i}] A_i$
\end{itemize}

\textbf{展开形式}：

\begin{align}
[\text{ad}_{V_i}](A_i) \dot{\theta}_i &= \text{ad}_{V_i}(A_i) \dot{\theta}_i \\
&= ([\text{ad}_{V_i}] A_i) \dot{\theta}_i \\
&= \begin{bmatrix}
[\omega_i] & 0 \\
[v_i] & [\omega_i]
\end{bmatrix} \begin{bmatrix}
\omega_{A_i} \\
v_{A_i}
\end{bmatrix} \dot{\theta}_i \\
&= \begin{bmatrix}
\omega_i \times \omega_{A_i} \\
v_i \times \omega_{A_i} + \omega_i \times v_{A_i}
\end{bmatrix} \dot{\theta}_i \\
&= \begin{bmatrix}
(\omega_i \times \omega_{A_i}) \dot{\theta}_i \\
(v_i \times \omega_{A_i} + \omega_i \times v_{A_i}) \dot{\theta}_i
\end{bmatrix}
\label{eq:notation2_expanded}
\end{align}

\section{数学等价性}

\subsection{基本等价关系}

从数学上看，这两种表示是\textbf{等价的}：

\begin{equation}
[\text{ad}_{V_i}] A_i \dot{\theta}_i = [\text{ad}_{V_i}](A_i) \dot{\theta}_i
\label{eq:equivalence}
\end{equation}

\textbf{证明}：

根据$\text{ad}$算子的定义：
\begin{equation}
\text{ad}_{V_i}(A_i) = [\text{ad}_{V_i}] A_i
\end{equation}

因此：
\begin{align}
[\text{ad}_{V_i}](A_i) \dot{\theta}_i &= \text{ad}_{V_i}(A_i) \dot{\theta}_i \\
&= ([\text{ad}_{V_i}] A_i) \dot{\theta}_i \\
&= [\text{ad}_{V_i}] A_i \dot{\theta}_i
\end{align}

\subsection{符号约定的说明}

虽然在数学上等价，但在不同的文献和推导中，可能会使用不同的符号约定：

\begin{enumerate}
\item \textbf{矩阵乘法表示}：$[\text{ad}_{V_i}] A_i \dot{\theta}_i$
\begin{itemize}
\item 强调这是矩阵与向量的乘法
\item 计算顺序明确：先矩阵乘法，再标量乘法
\end{itemize}

\item \textbf{算子作用表示}：$[\text{ad}_{V_i}](A_i) \dot{\theta}_i$
\begin{itemize}
\item 强调这是Lie括号运算
\item 更符合Lie代数的数学语言
\item 括号表示算子的作用
\end{itemize}
\end{enumerate}

\section{详细展开计算}

\subsection{完整的分量展开}

为了更清楚地看到区别，我们给出完整的分量展开。

\textbf{第一步：计算$[\text{ad}_{V_i}] A_i$}

设：
\begin{align}
\omega_i &= (\omega_{i,x}, \omega_{i,y}, \omega_{i,z})^T \\
v_i &= (v_{i,x}, v_{i,y}, v_{i,z})^T \\
\omega_{A_i} &= (\omega_{A_i,x}, \omega_{A_i,y}, \omega_{A_i,z})^T \\
v_{A_i} &= (v_{A_i,x}, v_{A_i,y}, v_{A_i,z})^T
\end{align}

则：
\begin{equation}
[\omega_i] = \begin{bmatrix}
0 & -\omega_{i,z} & \omega_{i,y} \\
\omega_{i,z} & 0 & -\omega_{i,x} \\
-\omega_{i,y} & \omega_{i,x} & 0
\end{bmatrix}
\end{equation}

\begin{equation}
[v_i] = \begin{bmatrix}
0 & -v_{i,z} & v_{i,y} \\
v_{i,z} & 0 & -v_{i,x} \\
-v_{i,y} & v_{i,x} & 0
\end{bmatrix}
\end{equation}

\textbf{第二步：计算$[\text{ad}_{V_i}] A_i$的分量}

上部分（角速度部分）：
\begin{align}
[\omega_i] \omega_{A_i} &= \omega_i \times \omega_{A_i} \\
&= \begin{bmatrix}
\omega_{i,y} \omega_{A_i,z} - \omega_{i,z} \omega_{A_i,y} \\
\omega_{i,z} \omega_{A_i,x} - \omega_{i,x} \omega_{A_i,z} \\
\omega_{i,x} \omega_{A_i,y} - \omega_{i,y} \omega_{A_i,x}
\end{bmatrix}
\end{align}

下部分（线速度部分）：
\begin{align}
[v_i] \omega_{A_i} + [\omega_i] v_{A_i} &= v_i \times \omega_{A_i} + \omega_i \times v_{A_i} \\
&= \begin{bmatrix}
v_{i,y} \omega_{A_i,z} - v_{i,z} \omega_{A_i,y} + \omega_{i,y} v_{A_i,z} - \omega_{i,z} v_{A_i,y} \\
v_{i,z} \omega_{A_i,x} - v_{i,x} \omega_{A_i,z} + \omega_{i,z} v_{A_i,x} - \omega_{i,x} v_{A_i,z} \\
v_{i,x} \omega_{A_i,y} - v_{i,y} \omega_{A_i,x} + \omega_{i,x} v_{A_i,y} - \omega_{i,y} v_{A_i,x}
\end{bmatrix}
\end{align}

\textbf{第三步：乘以$\dot{\theta}_i$}

\begin{align}
[\text{ad}_{V_i}] A_i \dot{\theta}_i &= \begin{bmatrix}
(\omega_i \times \omega_{A_i}) \dot{\theta}_i \\
(v_i \times \omega_{A_i} + \omega_i \times v_{A_i}) \dot{\theta}_i
\end{bmatrix} \\
&= \begin{bmatrix}
\begin{bmatrix}
(\omega_{i,y} \omega_{A_i,z} - \omega_{i,z} \omega_{A_i,y}) \dot{\theta}_i \\
(\omega_{i,z} \omega_{A_i,x} - \omega_{i,x} \omega_{A_i,z}) \dot{\theta}_i \\
(\omega_{i,x} \omega_{A_i,y} - \omega_{i,y} \omega_{A_i,x}) \dot{\theta}_i
\end{bmatrix} \\
\begin{bmatrix}
(v_{i,y} \omega_{A_i,z} - v_{i,z} \omega_{A_i,y} + \omega_{i,y} v_{A_i,z} - \omega_{i,z} v_{A_i,y}) \dot{\theta}_i \\
(v_{i,z} \omega_{A_i,x} - v_{i,x} \omega_{A_i,z} + \omega_{i,z} v_{A_i,x} - \omega_{i,x} v_{A_i,z}) \dot{\theta}_i \\
(v_{i,x} \omega_{A_i,y} - v_{i,y} \omega_{A_i,x} + \omega_{i,x} v_{A_i,y} - \omega_{i,y} v_{A_i,x}) \dot{\theta}_i
\end{bmatrix}
\end{bmatrix}
\end{align}

\section{在加速度递推中的应用}

\subsection{加速度递推公式中的使用}

在加速度递推公式中：
\begin{equation}
\dot{V}_i = \text{Ad}_{T_{i,i-1}}(\dot{V}_{i-1}) + [\text{ad}_{V_i}] A_i \dot{\theta}_i + A_i \ddot{\theta}_i
\end{equation}

这里使用$[\text{ad}_{V_i}] A_i \dot{\theta}_i$表示，等价于：
\begin{equation}
\dot{V}_i = \text{Ad}_{T_{i,i-1}}(\dot{V}_{i-1}) + [\text{ad}_{V_i}](A_i) \dot{\theta}_i + A_i \ddot{\theta}_i
\end{equation}

\textbf{物理意义}：

这一项$[\text{ad}_{V_i}] A_i \dot{\theta}_i$表示：
\begin{itemize}
\item 由于连杆$i$的速度$V_i$和关节$i$的螺旋轴$A_i$之间的相互作用
\item 产生的耦合加速度项
\item 这是科里奥利力和向心力的来源之一
\end{itemize}

\subsection{与速度递推公式的关系}

从速度递推公式：
\begin{equation}
V_i = \text{Ad}_{T_{i,i-1}}(V_{i-1}) + A_i \dot{\theta}_i
\end{equation}

对时间求导时，$A_i \dot{\theta}_i$项会产生：
\begin{align}
\frac{d}{dt}[A_i \dot{\theta}_i] &= A_i \ddot{\theta}_i + \dot{A}_i \dot{\theta}_i
\end{align}

但由于$A_i$在坐标系$\{i\}$中是常数，$\dot{A}_i = 0$。然而，$\frac{d}{dt}[\text{Ad}_{T_{i,i-1}}(V_{i-1})]$项会产生与$A_i \dot{\theta}_i$相关的耦合项，这就是$[\text{ad}_{V_i}] A_i \dot{\theta}_i$的来源。

\section{计算顺序的重要性}

\subsection{矩阵乘法的结合性}

虽然两种表示在数学上等价，但在实际计算中需要注意：

\begin{itemize}
\item $[\text{ad}_{V_i}] A_i \dot{\theta}_i$：先计算$[\text{ad}_{V_i}] A_i$（$6 \times 6$矩阵乘以$6 \times 1$向量，得到$6 \times 1$向量），再乘以标量$\dot{\theta}_i$
\item $[\text{ad}_{V_i}](A_i) \dot{\theta}_i$：先计算$[\text{ad}_{V_i}](A_i) = [\text{ad}_{V_i}] A_i$（得到$6 \times 1$向量），再乘以标量$\dot{\theta}_i$
\end{itemize}

\textbf{注意}：不能写成$[\text{ad}_{V_i}] (A_i \dot{\theta}_i)$，因为：
\begin{itemize}
\item $A_i \dot{\theta}_i$是$6 \times 1$向量
\item $[\text{ad}_{V_i}] (A_i \dot{\theta}_i)$是$6 \times 6$矩阵乘以$6 \times 1$向量，结果正确
\item 但这样写会改变计算的语义，因为$A_i \dot{\theta}_i$本身有特定的物理意义（关节运动产生的速度）
\end{itemize}

\section{总结}

\begin{enumerate}
\item \textbf{数学等价性}：$[\text{ad}_{V_i}] A_i \dot{\theta}_i = [\text{ad}_{V_i}](A_i) \dot{\theta}_i$

\item \textbf{表示方式}：
\begin{itemize}
\item $[\text{ad}_{V_i}] A_i \dot{\theta}_i$：强调矩阵乘法
\item $[\text{ad}_{V_i}](A_i) \dot{\theta}_i$：强调算子作用
\end{itemize}

\item \textbf{展开形式}：两者都展开为：
\begin{equation}
\begin{bmatrix}
(\omega_i \times \omega_{A_i}) \dot{\theta}_i \\
(v_i \times \omega_{A_i} + \omega_i \times v_{A_i}) \dot{\theta}_i
\end{bmatrix}
\end{equation}

\item \textbf{物理意义}：表示由于速度$V_i$和螺旋轴$A_i$之间的相互作用产生的耦合加速度项

\item \textbf{计算顺序}：先计算$[\text{ad}_{V_i}] A_i$（得到$6 \times 1$向量），再乘以标量$\dot{\theta}_i$
\end{enumerate}

\section{参考文献}

\begin{itemize}
\item Modern Robotics: Mechanics, Planning, and Control, Chapter 8
\item 牛顿-欧拉逆向动力学详细推导：\texttt{newton\_euler.tex}
\end{itemize}

\end{document}



